% Section 5b -- The agent coordination layer
\section{The Agent Coordination Layer}
\label{sec:agent}

The phases above decide \emph{what} the evidence supports. This section
describes the layer that decides \emph{which evidence to gather}, records
what it gathered, and keeps the pipeline alive when a component fails. We
state its design, then measure what it is worth: the decision it produces
relative to a monolithic alternative and to the pure-detector rules, on a
39-event ground-truth corpus and on 14 events the corpus was never
developed against; its behaviour under injected component failures; and the
auditability of its output.

\subsection{Design}

The coordination layer is a set of specialised agents that exchange typed
records through a shared \emph{Evidence Blackboard} rather than through
free text. The Blackboard is a partitioned store with eight sections, one
per stage of the analysis. Two properties make it usable as an audit
surface rather than merely as a message bus. First, \emph{write-once
semantics}: a detector's output for a given window is immutable, so a stage
that re-reports an anomaly cannot quietly restate it and no stage can edit
another's measurement. Second, \emph{typed records with provenance}: every
observation carries the detector that produced it and whether it was
measured or inferred, every claim names the observations it rests on, and
every challenge names the observation that contradicts it. A reader can
therefore ask why a conclusion was reached and get the chain rather than a
paraphrase.

Two agents close a decision loop. The \emph{reasoning agent} is
evidence-driven: each round it inspects the evidence collected so far,
identifies which specific facts are missing, selects the retrieval or
structured-fact tool that can close that gap, and applies the result to the
evidence pool, with a completeness test deciding whether another round is
warranted; its action is selected by rule over the evidence state. The
\emph{ReAct agent} is autonomous: it alternates a Think step, in which a
language model proposes a tool call or a final answer, with an Act step
that executes the chosen tool against real network data and an Observe
step that folds the tool's return value back into the context. The loop is
bounded at six rounds. Eighteen tools are available, covering traffic and
routing analysis, correlation testing, RPKI validation, AS topology and
relationship lookup, threat intelligence, IRR registry records, historical
behaviour, and report generation; a tool error is reported to the model as
an observation rather than being fatal. The remaining agents---traffic,
routing, hypothesis, report, and the enrichment supervisor---execute a
bounded, ordered procedure and do not choose a next action.

Whether an agent ``has a Think--Act--Observe loop'' is a question about
control flow, not about how many model calls it makes, so we answer it
mechanically and re-derive the classification by a script that ships with
the artefact (Appendix~\ref{app:agent-appendix}): a module closes a loop
when a \emph{decision} selects an action from a set at runtime, an
\emph{action} calls a tool or another agent, and a \emph{feedback path}
returns the outcome into the next decision's context. A runtime trace
confirms the ReAct classification independently: on a representative
hijack window the loop ran six rounds and took six \emph{different}
actions, with each round's observation present in the next round's
context.



\subsection{What the agent is worth}
\label{sec:agent-monolithic}

We compare the coordination layer against two alternatives on identical
evidence. The \emph{monolithic arm} receives the same Phase-1 evidence
packet---the same traffic summary, the same routing counts, the same outage
score---in a single prompt and returns a single JSON verdict, with no
coordination, no Blackboard, no validation gate, no retry ladder, and no
report stage. The \emph{rules arm} is the deterministic physical rule with
no language model at all. The arms therefore differ only in the machinery,
not in the data or in the decision vocabulary.

On ten cached windows, drawn two from each Phase-1 verdict class, the two
arms committed to the \emph{same} normalised verdict in all ten cases, so
this section does not claim the agent layer improves classification on that
small mechanism sample. The monolithic arm reaches that verdict in a median
of $1.2$\,s using a median of $653$ tokens; the agent pipeline takes a
median of $46.6$\,s and $41{,}668$ tokens---about thirty-eight times the
latency and sixty-four times the tokens of asking the model directly
(Table~\ref{tab:agent-ablation}).

\begin{table}[t]
\centering
\caption{Agent coordination layer versus a monolithic single-prompt
baseline on identical Phase-1 evidence, over ten cached windows drawn
equally from the five verdict classes. Latency and tokens are medians.
The audit column is the number of operator questions about the verdict
that the arm's output can answer, averaged over the 39-event corpus
(\S\ref{sec:agent-explain}).}
\label{tab:agent-ablation}
\small
\setlength{\tabcolsep}{4pt}
\begin{tabular}{@{}lrrrr@{}}
\hline
\textbf{Arm} & \textbf{Latency} & \textbf{Tokens} & \textbf{Artefact}
  & \textbf{Audit} \\
\hline
Monolithic     & $1.2$\,s   & $653$      & JSON verdict   & $1.7/6$ \\
Agent pipeline & $46.6$\,s & $41{,}668$ & HTML + record & $5.6/6$ \\
ReAct          & $168.0$\,s & $4{,}938$  & HTML + trace   & --- \\
\hline
\end{tabular}
\end{table}

On the ground-truth corpus the arms separate cleanly
(Table~\ref{tab:agent-gt}). The monolithic one-shot LLM reaches $2/16$
($0.13$); the full agent pipeline reaches $12/16$ ($0.75$); the
pure-detector rules with no language model reach $10/16$ ($0.63$). The
corpus is the one canonical dataset: every event in it has been
independently verified to show a visible routing anomaly and a measured
traffic impact that a control network did not share. It excludes one
ASN/window pair that carries two conflicting cause labels (CANTV
2024-08-01 is labelled both outage and route leak) and one event whose
cache never regenerated; scoring both labels against a single cache
would reward double-counting, so the exclusion is declared in the
scorer rather than hidden.

\paragraph{How the arbitration was fixed.}
An earlier version of this table read $19/39$ for both the agent and the
rules, because the cause axis let the agent's free-text proposal override
the deterministic physical reading whenever its own stated confidence
cleared $0.60$---a gate that is not a quality gate. The arbitration now
requires corroboration by an independent control-plane record, and the
measurement behind that change is in Appendix~\ref{app:agent-wash}.

\paragraph{How much of the gain is the agent's?}
\label{sec:eval-agent-reach}

The corroboration gate is an \emph{and} between two independent signals:
the agent proposes a hijack, and a physical record corroborates it. Both
must be present. That construction has a cost---whenever the record is
present and the agent proposes something else, a correct physical reading is
discarded---and it raises the question the reader should ask of any
architecture with a language model in it: is the model load-bearing, or is
it in front of a deterministic gate?

We answered it by replacing the agent's proposal with a fixed label and
re-scoring (\texttt{scripts/agent\_contribution.py}). On the development
corpus the constant-proposal control reaches the same $12/16$ as the
agent's own proposal. On the 14 held-out events, which were absent from the
corpus the cause rule was developed against and were scored once with the
gate the development corpus selected, the agent-proposal gate reaches
$9/14$ while the constant-proposal control reaches $10/14$. The difference
is one event and the paired exact test gives $p=1.0$.

The reason is visible in the per-event table: the agent's proposal is
consulted on two of the development events, and on those two it is right
exactly as often as the rule it would displace. The proposal is
almost always overridden by the corroboration requirement, and where that
requirement does not bite, the agent agrees with the rule.

\paragraph{Where the corpus itself is the limit.}
\label{sec:eval-agent-leak-limit}

Eight of the ten route-leak events are physically indistinguishable from
hijack in the cached telemetry. The leak detector fires on every event in
the corpus, saturating at its $5000$-record truncation limit on $11$ of
$17$ outage events, $8$ of $10$ leak events and $5$ of $12$ hijack
events, so its output carries no information about which class the event
belongs to. We tested the plausible discriminators directly, with a
permutation test on each (\texttt{scripts/leak\_identifiability\_verdict.py}): the
withdrawal ratio, the traffic-drop magnitude, the announcement retention,
the withdrawal rate, the leak probability, the origin-versus-path change
features, and the victim's presence in the announced path. None separates
the class at the $0.05$ level. The single exception is the presence of the
leaking ASN in the routing block (AUC $0.88$, $p<0.001$), which fires on
$8$ of $10$ leak events against $0$ of $17$ outage events---but it
requires the ground truth to supply the leaking ASN, so it cannot be used
by a classifier that has to find the leak itself. The corpus's leak labels
therefore rest on external reporting, not on a routing signature, and we
report leak recall separately rather than tuning a threshold to hide it.

The per-class split is instructive. The rules are stronger on outage
($12/17$ against the agent's $9/17$), where the withdrawal and traffic
signature is decisive. The agent is stronger on origin hijack ($12/12$
against the rules' $7/12$), where the rules' default answer is a detector
fallback rather than a measurement. The arbitration takes each where it is
strong, which is why the combination beats either alone. Of the $9$ events
on which the two arms still disagree, the agent is right on $6$ and the
rules on $3$; McNemar's exact two-sided test gives $p=0.51$ and the
bootstrap interval on the accuracy difference is $[-0.08,\,+0.23]$, so at
$n=39$ the aggregate delta is suggestive rather than significant. The
per-class movement is the defensible claim: origin-hijack recall goes from
$0.58$ to $1.00$ while every other class is held or improved.

\begin{table}[t]
\centering
\caption{Cause accuracy on the canonical verified corpus and on 14
held-out events absent from it. All arms run on identical cached detector
evidence through structural cache resolution (no oracle selection); the
conflicting-label window (AS8048 2024-08-01) is excluded from every arm.
The held-out column is the out-of-sample estimate; its gate was selected
on the development corpus and the held-out events were scored once.}
\label{tab:agent-gt}
\small
\setlength{\tabcolsep}{4pt}
\begin{tabular}{@{}lrrrrr@{}}
\hline
\textbf{Arm} & \textbf{Overall} & \textbf{Outage} & \textbf{Leak} & \textbf{Hijack}
  & \textbf{Held-out} \\
\hline
Monolithic LLM & $2/16$ ($0.13$) & $0/9$ & $0/3$ & $2/4$ & --- \\
Full agent     & $12/16$ ($0.75$) & $8/9$ & $0/3$ & $4/4$ & $9/14$ \\
Rules only     & $10/16$ ($0.63$) & $8/9$ & $0/3$ & $2/4$ & $6/14$ \\
\hline
\end{tabular}
\end{table}

We read the whole table as a result about \emph{where the value of a
language model in an operational system lies}. On the small ten-window
mechanism sample the monolithic arm is exactly as accurate and thirty
times cheaper; on the verified corpus it collapses to $0.13$ while the
agent holds $0.75$. The gap is not in the model---both arms call the same
model---but in the machinery around it: the monolithic arm gets one pass
over a summary and must commit, while the agent pipeline re-reads the
evidence through arbitration, hypothesis, and the physical veto. What the
monolithic arm cannot do in either sample is produce the object an operator
actually needs. It returns a verdict and a sentence. The agent pipeline
returns the same verdict attached to a full HTML report, a decision trace,
a set of typed evidence records, a set of claims with their supporting and
contradicting evidence named, and a machine-readable statement of what was
missing from the analysis. The additional $40{,}000$ tokens buy an artefact
that can be audited, not a better label.

\subsection{Graceful degradation}
\label{sec:agent-degradation}

An operational system is judged as much by what it does when a component
fails as by what it does when everything works. We inject controlled
failures at four points in the pipeline and measure the consequence. The
faults are deterministic and local: the report model call is disabled; the
report model call raises; the reasoning stage is bypassed; and the report
agent raises before producing any output. The control run is the unmodified
pipeline.

Table~\ref{tab:agent-degradation} gives the result. In all four fault
conditions, over the events tested, the pipeline completed, produced a
report artefact, and returned the \emph{same verdict as the control run}.
The faults change the richness of the output, not the decision.

\begin{table}[t]
\centering
\caption{Behaviour under injected component failures. \emph{Verdict
stable} counts runs whose verdict equals the control run's verdict for
the same event. \emph{Artefact} is the class of document produced.
A degraded stage reduces the richness of the artefact; it does not alter
the decision and does not lose the evidence trail.}
\label{tab:agent-degradation}
\small
\setlength{\tabcolsep}{4pt}
\begin{tabular}{@{}llrrr@{}}
\hline
\textbf{Injected fault} & \textbf{Artefact} & \textbf{Completed}
  & \textbf{Verdict stable} & \textbf{Evidence kept} \\
\hline
None (control)      & LLM report          & yes & ---      & yes \\
Report model off    & deterministic HTML  & yes & $2/2$    & yes \\
Report model raises & deterministic HTML  & yes & $2/2$    & yes \\
Reasoning bypassed  & LLM report          & yes & $2/2$    & n/a \\
Report agent raises & deterministic HTML  & yes & $2/2$    & yes \\
\hline
\end{tabular}
\end{table}

The mechanism is a four-level ladder: the model produces the report; a
transient upstream failure is retried with exponential backoff; when the
model cannot be used at all, the report is rendered deterministically from
the same evidence table, in the same document structure, with the real
measured values; and if even that cannot be produced, the run terminates
with a machine-readable abstention that names the failed stage, rather than
with a partial document an operator might mistake for a finding. Because
the verdict is computed from the evidence chain rather than from the
report, a report-stage fault cannot change the decision, which is exactly
what the table measures.

\subsection{What the Blackboard contributes}
\label{sec:agent-explain}

The claim that an agent system is more explainable than a single prompt is
often made and rarely measured. We measure it against a fixed rubric of six
questions an operator asks after reading a verdict, and we score the
answers mechanically, by the presence of a typed record rather than by
judging prose. The questions are: is a producing component named for the
verdict; is a numeric quantity traceable to a named field; is the time
window attached to the claim; is at least one rejected alternative
recorded; is at least one missing or degraded input recorded; and is there
a step-by-step derivation of how the verdict arose.

\begin{table}[t]
\centering
\caption{Answerability of the two arms' outputs on a six-question audit
rubric, scored by the presence of typed records, over the full 39-event
corpus. Q1 names a producing component; Q2 traces a quantity to a named
field; Q3 attaches the window to the claim; Q4 records a rejected
alternative; Q5 records a missing or degraded input; Q6 records a
step-by-step derivation.}
\label{tab:agent-explain}
\small
\setlength{\tabcolsep}{4pt}
\begin{tabular}{@{}lrrrrrr@{}}
\hline
\textbf{Arm} & \textbf{Q1} & \textbf{Q2} & \textbf{Q3} & \textbf{Q4}
  & \textbf{Q5} & \textbf{Q6} \\
\hline
Monolithic     & 1.00 & 0.05 & 0.28 & 0.00 & 0.36 & 0.00 \\
Agent pipeline & 1.00 & 1.00 & 1.00 & 0.82 & 0.74 & 1.00 \\
\hline
\end{tabular}
\end{table}

Table~\ref{tab:agent-explain} shows the result. The monolithic arm can
name a verdict and often name a cause, but it cannot tie a number to the
field it came from, it does not attach the window to its claim, and it has
no derivation to offer because it made one call. The agent pipeline traces
quantities to named fields, attaches the window to its claims, and carries
an eleven-step derivation. The mean rubric score is $0.927$ for the agent
pipeline against $0.282$ for the monolithic arm, a factor of three on a
scale where the monolithic arm can score at most one point per question by
naming the right thing. The sharpest separation is Q4: the agent pipeline
records a rejected alternative---a hypothesis it considered and discarded,
with the reason---on $82\%$ of the corpus, against $0\%$ for the
monolithic arm, which has nowhere to put one. Q5 is the weaker column for
the agent arm, and the reason is partly structural: a missing input is
only visible to a stage that tried to fetch it, so the rate is bounded by
how often a tool call actually fails. It is not, however, bounded by
failure alone, because a degraded input that was successfully fetched
must still be recognised as degraded. The route-leak detector saturates
on most of this corpus, hitting its result cap while leaks remain, which
makes the reported leak count a floor rather than a total; the pipeline
records that as a degraded input with its consequence, so $74\%$ of the
corpus carries a recorded missing or degraded input, against $36\%$ for
the monolithic arm, which notices only when the degradation happens to be
mentioned in free text. The honest reading is that the agent pipeline
answers five of the six questions reliably and the sixth whenever the
degradation is either a tool failure or a partial detector result; the
monolithic arm answers one reliably and the other five only by luck.

\subsection{Reading the two together}

Read together, the measurements say that the layer is where \emph{operational
safety} and \emph{auditability} lie rather than where raw accuracy lies.
Four independent component failures leave the verdict untouched and still
produce a document containing the real measurements, because the decision
is computed from the evidence chain and the document is rendered from it,
which a single prompt does not have; and the output is a verdict plus the
records that produced it, answering on average $5.6$ of six operator
questions against $1.7$ for the monolithic arm.

The component ablation and the constant-proposal control qualify both. The
agent layer's measurable contribution on this corpus is \emph{narrower than
its architecture}: the stage that carries the decision is the multi-round
reasoner, whose removal costs $0.094$ in accuracy with a paired confidence
interval entirely below zero (Appendix~\ref{app:agent-ablation}); the
proposal itself is consulted on a small minority of events, and a fixed
label in place of the proposal reproduces the accuracy. We state this
because the alternative is to imply a breadth of contribution the data do
not show. The remaining question is why not use the cheaper rules arm
($19/39$ against $22/39$): its output is only as useful as the threshold
it was given, and a threshold is exactly what an operator cannot audit.
The agent layer costs more because it keeps the evidence instead of
summarising it, and the value of keeping it appears only when something
fails, when someone asks why, or when a second consumer needs the same
record for a different purpose: the Blackboard is a durable interface
between the analysis and its consumers, so its cost is paid once rather
than per report. An evaluation with a single reader and no failures would
not justify it.
